\subsection{Countable targets for reverse commutator towers}

\begin{claim}{A.21}{Proved}{Reverse commutator towers still have countable images}
Use the convention $[u,v]=uvu^{-1}v^{-1}$. For a group $H$ and
sequences $a_n,b_n\in H$, define
\[
 E(H;a,b)=\left\langle H,x_n\ (n\geq0)\ \middle|\
 x_n=[x_{n+1},a_n]b_n\ (n\geq0)\right\rangle.
\]
The canonical map $H\to E(H;a,b)$ is injective. An aspherical CW
model $X$ for $H$ extends by the displayed one- and two-cells to
an aspherical model for $E(H;a,b)$; the inclusion induces an
isomorphism on ordinary integral homology. Consequently it preserves
an actual geometric-dimension bound of two and preserves acyclicity.
The larger group is normally generated by the image of $H$.

Nevertheless, for every nontrivial countable torsion-free group $H$
there is a fixed countable torsion-free group $Q$ containing $H$
such that every set-sized transfinite iteration of these extensions,
starting with $H$ and using group unions at limit stages, admits a
homomorphism to $Q$ restricting to the specified inclusion $H\to Q$.
Its image is a nontrivial countable quotient. Thus arbitrarily long
iterations of these particular homology-preserving two-dimensional
extensions cannot remove all countable images.
\end{claim}
\begin{proof}
Start with $X\vee\bigvee_{n\geq0}S^1_{x_n}$ and adjoin the two-cells
in the order $n=0,1,\ldots$. After adjoining the first $N+1$ cells,
each generator $x_n$, for $0\leq n\leq N$, can be eliminated, beginning
with $x_0$, using its unique occurrence as the left side of the
corresponding relation. More precisely, successive substitution
identifies that group's presentation with
\[
 H*F(x_{N+1},x_{N+2},\ldots).
\]
These substitutions are elementary presentation homotopy equivalences:
the eliminated one-cell and its defining two-cell form a cancelling
pair after the corresponding elementary two-cell slides. They retain
$X$ and its homotopy type. Hence every finite initial stage is
aspherical and its map from $H$ is injective. Each sphere of dimension at least two in the full
CW union has image in finitely many cells and therefore lies in one
such stage, where it is nullhomotopic. The union is aspherical. A word
from $H$ that becomes trivial in the union is already trivial at a
finite stage, so the indicated map from $H$ remains injective.

The relative ordinary cellular chain complex has one generator $e_n$
in degree one and one generator $r_n$ in degree two, with
$\partial r_n=e_n$ after choosing orientations: the exponent sum of
$x_{n+1}$ in its commutator is zero. There are no other relative
cells. This chain complex is contractible, proving the homology
isomorphism. Killing $H$ in the presentation makes every relation
$x_n=1$, proving normal generation. The same asphericity,
injectivity and homology conclusions persist at transfinite unions,
by the compact finite-cell argument and continuity of cellular
homology.

We construct the fixed target independently. Set $Q_0=H$. For each
integer $k\geq0$, adjoin to $Q_k$ a stable letter $t_{u,v}$ for every
ordered pair of nonidentity elements $u,v\in Q_k$, with relation
$t_{u,v}^{-1}ut_{u,v}=v$. The associated cyclic subgroups are infinite,
so these are genuine HNN extensions. Normal forms give injectivity
of $Q_k\to Q_{k+1}$. The action on the Bass--Serre tree shows that a
finite-order element is conjugate into the torsion-free vertex group;
hence $Q_{k+1}$ is torsion-free. All groups in this countable sequence
are countable. Their union $Q$ is countable, torsion-free, contains
$H$, and has exactly two conjugacy classes: the identity class and
the class of all nonidentity elements. Indeed, two such elements
belong to a common $Q_k$ and become conjugate in $Q_{k+1}$.

For $a,b\in Q$ with $a\ne1$, consider
\[
 f_{a,b}:Q\longrightarrow Q,\qquad
 f_{a,b}(y)=[y,a]b= y a y^{-1}a^{-1}b.
\]
All nonidentity elements are conjugate to $a$, so
\[
 \operatorname{im} f_{a,b}=Q\setminus\{a^{-1}b\}.
\]
Every nonempty fibre is a right coset of $C_Q(a)$ and is infinite,
since that centralizer contains the infinite cyclic group $\langle
 a\rangle$. If $a=1$, the map $f_{a,b}$ is the constant map with
value $b$.

We now solve every system $x_n=f_{a_n,b_n}(x_{n+1})$ within $Q$.
First suppose all $a_n\ne1$. Choose $x_0$ away from the missing
point of $f_{a_0,b_0}$. Given $x_n$ in the image of $f_{a_n,b_n}$,
choose a preimage $x_{n+1}$ that avoids the missing point of
$f_{a_{n+1},b_{n+1}}$. The relevant fibre is infinite and only one
value is excluded, so this recursion continues indefinitely and
satisfies every equation.

An index $r$ with $a_r=1$ instead specifies $x_r=b_r$ and places no
condition on $x_{r+1}$. Between consecutive such indices, and before
the first, solve the finite block backwards from its prescribed
last value. If there are infinitely many such indices these finite
blocks cover all equations. If there are finitely many, the last
remaining infinite block has no zero $a_n$ and is solved by the
preceding recursion. This includes the case of no zero $a_n$.

Given a homomorphism $\phi:G\to Q$ from any current stage, apply
this solution procedure to the coefficient sequences
$\phi(a_n),\phi(b_n)$. The resulting values of the $x_n$ extend
$\phi$ to $E(G;a,b)\to Q$, by its presentation. At limit stages take
the union of the compatible homomorphisms. Transfinite recursion
therefore produces the asserted map from every iteration to the
same $Q$, and its restriction to $H$ is injective. Its image is
countable and nontrivial, as required.
\end{proof}
